% Standalone problem document, generated from the adjacent README.md.
% Compile with XeLaTeX. Mathematical content is maintained in Markdown.
\documentclass[11pt,a4paper]{article}
\usepackage[fontsize=10.5pt]{scrextend}
\usepackage[margin=22mm,headheight=15pt,headsep=9mm,footskip=11mm]{geometry}
\usepackage{fontspec}
\setmainfont{texgyrepagella}[Extension=.otf,UprightFont=*-regular,BoldFont=*-bold,ItalicFont=*-italic,BoldItalicFont=*-bolditalic]
\setsansfont{texgyreheros}[Extension=.otf,UprightFont=*-regular,BoldFont=*-bold,ItalicFont=*-italic,BoldItalicFont=*-bolditalic]
\setmonofont{texgyrecursor}[Extension=.otf,UprightFont=*-regular,BoldFont=*-bold,ItalicFont=*-italic,BoldItalicFont=*-bolditalic,Scale=MatchLowercase]
\usepackage{amsmath,amssymb}
\usepackage{unicode-math}
\setmathfont{texgyrepagella-math.otf}
\usepackage{xcolor,microtype,enumitem,needspace,fancyhdr}
\usepackage{bookmark}
\definecolor{ink}{HTML}{17344B}
\definecolor{accent}{HTML}{197D80}
\definecolor{muted}{HTML}{596773}
\hypersetup{colorlinks=true,linkcolor=accent,urlcolor=accent,pdftitle={MF-15 - Critical
exponent for generalized doubly nonnegative
matrices},pdfauthor={Open Problems in Numerical Linear Algebra}}
\urlstyle{same}
\setlength{\parindent}{0pt}
\setlength{\parskip}{5pt plus 1pt minus 1pt}
\linespread{1.04}
\setlength{\emergencystretch}{3em}
\widowpenalty=10000
\clubpenalty=10000
\displaywidowpenalty=10000
\allowdisplaybreaks[1]
\setlist{nosep,leftmargin=1.5em}
\providecommand{\tightlist}{\setlength{\itemsep}{0pt}\setlength{\parskip}{0pt}}
\providecommand{\pandocbounded}[1]{#1}
\pagestyle{fancy}
\fancyhf{}
\fancyhead[L]{\sffamily\footnotesize\color{muted}OPEN PROBLEMS IN NUMERICAL LINEAR ALGEBRA}
\fancyhead[R]{\sffamily\bfseries\footnotesize\color{accent}MF-15}
\fancyfoot[L]{\parbox[b]{0.9\textwidth}{\sffamily\fontsize{8}{10}\selectfont\color{muted}Editorial ratings; a bounded search cannot certify that a problem remains unsolved.\\Literature check: 2026-09-10}}
\fancyfoot[R]{\sffamily\footnotesize\color{muted}\thepage}
\renewcommand{\headrulewidth}{0.3pt}
\renewcommand{\headrule}{\hbox to\headwidth{\color{accent}\leaders\hrule height \headrulewidth\hfill}}
\makeatletter
\renewcommand{\section}{\Needspace{5\baselineskip}\@startsection{section}{1}{0pt}{12pt plus 2pt minus 2pt}{4pt}{\sffamily\bfseries\large\color{ink}}}
\renewcommand{\subsection}{\Needspace{4\baselineskip}\@startsection{subsection}{2}{0pt}{9pt plus 1pt minus 1pt}{3pt}{\sffamily\bfseries\normalsize\color{ink}}}
\makeatother
\setcounter{secnumdepth}{0}
\begin{document}
{\sffamily\small\color{muted}Matrix functions and stability\par}
\vspace{3pt}
{\raggedright\hyphenpenalty=10000\exhyphenpenalty=10000\sffamily\bfseries\fontsize{21}{24}\selectfont\color{ink}Critical
exponent for generalized doubly nonnegative matrices\par}
\vspace{7pt}
{\sffamily\small\textbf{Difficulty:} challenging\quad\textbf{Importance:} interesting
to the community\par}
{\sffamily\small\bfseries\color{accent}Status: Partially resolved\par}
\vspace{4pt}
{\color{accent}\hrule height 0.8pt}
\vspace{6pt}
\textbf{Topic:} Positivity of conventional matrix powers.

\textbf{Rating rationale:} Challenging because nonsymmetric spectral
powers require a sharp dimension-uniform positivity threshold; community
impact concerns matrix functions and positive matrix dynamics.

\subsection{Problem statement}\label{problem-statement}

For \(n\ge3\), let \(\mathcal G_n\) consist of the matrices
\(A\in\mathbb R^{n\times n}\) that are entrywise nonnegative,
diagonalizable, and have only nonnegative real eigenvalues. Symmetry is
not assumed. If

\[
A=S\operatorname{diag}(\lambda_1,\ldots,\lambda_n)S^{-1},
\]

define the conventional matrix power, for \(t>0\), by

\[
A^t=S\operatorname{diag}(\lambda_1^t,\ldots,\lambda_n^t)S^{-1},
\qquad 0^t=0.
\]

Let

\[
\mathrm{CE}_n=\inf\bigl\{c\ge0:
A^t\text{ is entrywise nonnegative for every }A\in\mathcal G_n
\text{ and every real }t>c\bigr\}.
\]

Is \(\mathrm{CE}_n=2(n-2)\) for every integer \(n\ge3\)?

\subsection{Why it matters}\label{why-it-matters}

This asks for the exact dimension-dependent threshold beyond which
spectral matrix powers preserve entrywise positivity in a nonsymmetric
class, relevant to matrix functions and positive dynamics.

\newpage
\subsection{References}\label{references}

\begin{itemize}
\tightlist
\item
  X. Han, C. R. Johnson, and P. Paparella,
  \href{https://arxiv.org/abs/1407.7059}{The critical exponent for
  generalized doubly nonnegative matrices}, \emph{Linear and Multilinear
  Algebra} 65 (2017), 1035--1044, §1 for definitions, Theorem 3.3 for an
  upper bound, Corollary 4.5 for \(n=3\), and Question 4.7 for the
  displayed formula. The conjectural \(n=4\) case is Conjecture 4.6.
\item
  D. Guillot, A. Khare, and B. Rajaratnam,
  \href{https://arxiv.org/abs/1303.4701}{The critical exponent
  conjecture for powers of doubly nonnegative matrices}, \emph{Linear
  Algebra and its Applications} 439 (2013), 2422--2427, abstract and
  main theorem: resolution of the symmetric case.
\end{itemize}

\subsection{Status check}\label{status-check}

On 2026-09-08, checked the latest listed arXiv version and searched
``generalized doubly nonnegative'', ``critical exponent'', ``2(n-2)'',
``conjecture'', and 2025--2026. No resolution of Question 4.7 was
located. The source proves existence of a finite threshold and
\(\mathrm{CE}_3=2\). The solved symmetric doubly nonnegative problem has
threshold \(n-2\) and concerns a smaller class; results for entrywise
(Hadamard) powers concern a different operation. The special case
\(n=4\) and related integrality questions are included within this
single target, not counted separately. No recent explicit reaffirmation
was found.

\subsection{Audit --- 2026-09-10}\label{audit-2026-09-10}

Rechecked
\href{https://arxiv.org/pdf/1407.7059}{Han--Johnson--Paparella,
Corollary 4.5 and Question 4.7}. The \(n=3\) equality is proved; the
all-dimensions formula remains posed. Generalized-DN critical-exponent
searches found no later resolution. The
\href{https://doi.org/10.1080/03081087.2016.1223009}{published record}
confirms the 2017 article; symmetric and entrywise-power theorems do not
settle it.
\end{document}
